%! Characteristics of Polynomial Functions II — Unit 3, Lesson 3.5 — Student Packet Cover Sheet
% UNITS 3+ PACKET-FIRST shape (LESSON_SHAPE.md §0a, 2026-11-14). 12pt; \coverbanner measures the
% title block and sizes the band. THREE rows (no extra credit, 2026-11-17) (a contents list, no Score column): Warm-Up, Guided
% Notes & Practice, Homework (DeltaMath online — not in this packet).
% Homework is due the first class after two study halls — never "next class".
% Keep in Mind = the Closing Practice's Watch out.
\documentclass[12pt]{article}
\usepackage{algebra2-article}
\usepackage{algebra2-boxes}

\begin{document}

\coverbanner{Unit 3 \quad Polynomial Functions}{Lesson 3.5 \quad Characteristics of Polynomial Functions II}

\namedateperiod
\vspace{0.18in}

\boxguard[14]
\begin{learningtargetbox}
\small
\begin{enumerate}[leftmargin=1.5em, itemsep=5pt]
  \item \ldots{} find each \textbf{zero} of a polynomial in factored form and its
        \textbf{multiplicity}, and tell whether the graph \textbf{crosses} (odd multiplicity) or
        \textbf{touches} (even multiplicity) the $x$-axis there --- and read the zeros off a graph.
  \item \ldots{} count a graph's \textbf{turning points} (at most $n-1$ for degree $n$) and name its
        \textbf{relative maximum} and \textbf{minimum} values and any \textbf{absolute maximum} or
        \textbf{minimum}.
  \item \ldots{} write the intervals where a polynomial is \textbf{increasing},
        \textbf{decreasing}, \textbf{positive}, and \textbf{negative} in \textbf{interval
        notation} using $x$-values, and state its \textbf{domain} and \textbf{range}.
\end{enumerate}
\end{learningtargetbox}

\vspace{0.14in}

\boxguard[14]
\begin{tocbox}
\small
\renewcommand{\arraystretch}{1.5}
\begin{tabularx}{\linewidth}{@{} c l X @{}}
\rowcolor{forestbg}
\textbf{\#} & \textbf{Component} & \textbf{Description} \\
\midrule
1 & Warm-Up & Degree, end behavior and symmetry from Lesson 3.4, and the zeros of a factored
              polynomial \\
2 & Guided Notes \& Practice & Three rounds --- zeros and multiplicity, turning points and
              extrema, intervals and domain/range --- each a definition, an example, and a
              one-to-three-star practice ladder; then the Closing Practice, scored out of 4 \\
3 & Homework & Assigned on DeltaMath: \textbf{A2: Lesson 3.5: Zeros, Extrema \& Intervals} --- complete it
              online; it is not in this packet. \textbf{Scored 1 to 5, due the first class after
              two study halls.} \\
\end{tabularx}
\end{tocbox}

\vspace{0.14in}

\boxguard[8]
\begin{remindbox}
\small
A \textbf{zero} of \textbf{odd multiplicity} \emph{crosses} the $x$-axis; one of \textbf{even
multiplicity} \emph{touches} it and turns back. A polynomial of degree $n$ has \emph{at most}
$n-1$ \textbf{turning points}. A \textbf{relative} maximum or minimum is the highest or lowest
point \emph{nearby}; an \textbf{absolute} one is the highest or lowest of the \emph{whole}
function --- an odd-degree polynomial has none, so not every maximum is ``the'' maximum.
\textbf{Watch out:} \textbf{increasing}, \textbf{decreasing}, \textbf{positive}, and
\textbf{negative} intervals are written with \textbf{$x$-values}, read along the $x$-axis from
turning point to turning point (or zero to zero) --- never with a turning point's $y$-value.
\end{remindbox}

\end{document}
